% GENERATED FILE. Edit canonical lessons, then run npm run generate:books.
% canonical-source-hash: fnv1a64:35e11539b70daa41
% canonical-lessons: ZH-W10-yan, ZH-W10-xie, ZH-C10-xiexie, ZH-C10-buxie, ZH-C10-practice

\chapter{One Character, Two Courtesies}
\label{ch:zh-xiexie}
\hlchaptermodality{10}

\begin{chapteropening}
\textbf{By the end of this chapter:} \emph{I can write the speech radical \zh{讠} and the twelve-stroke \zh{谢}, thank someone with \zh{谢谢}, answer thanks with \zh{不谢}, and predict from the rule alone that \zh{不} rises before \zh{谢} because \zh{谢} is a falling tone.}

A whole courtesy exchange bought with one new character -- doubled for the thanks, negated for the reply -- and the first place the reader predicts a tone change from a rule rather than being told how a word sounds.
\end{chapteropening}

\section[yán]{\zh{讠} — the speech radical}

\label{lesson:ZH-W10-yan}



\begin{quote}
Write \zh{亻}. Two strokes, squeezed to the left edge of a block.
\end{quote}

\subsection*{Script — the shape}
\begin{quote}
\zh{讠}
\end{quote}

\textbf{Two strokes}, one pen lift, and it is not a character. Like \zh{亻}, it is a \textbf{component} — a piece that lives at the left edge of other characters.

1. a \textbf{dot} at the top, pressed down and to the right 2. a short \textbf{horizontal}, which \textbf{turns down} and then \textbf{finishes rising} to the upper right — all one stroke, no lift

That second stroke does three things without stopping. You have handled worse: the frame of \zh{再}, the leg of \zh{见}, the turn in \zh{么}.

\textbf{\zh{讠} means speech.} It is the squeezed form of \zh{言}, the character for words and speaking, and it does the same job \zh{亻} does: it tells you what the character is \emph{about} before you know how to say it.

\zh{讠} joins \zh{亻} as a component that carries \textbf{meaning} rather than sound, and the pair is worth holding together:

\begin{itemize}
\raggedright
  \item \textbf{\zh{亻}} at the left edge $\to$ the character is about a \textbf{person}
  \item \textbf{\zh{讠}} at the left edge $\to$ the character is about \textbf{speech}
\end{itemize}

You will meet \zh{讠} often. Every time you do, you can guess something true about the word before anyone tells you.

\subsection*{Your turn}
Write these out:

\begin{itemize}
\raggedright
  \item \zh{讠} — dot, then one unbroken turning stroke
  \item \zh{亻}, then \zh{讠}, side by side — two meaning-carrying components
\end{itemize}

\begin{tcolorbox}[breakable,colback=teal!4,colframe=teal!35!black,title={Before you move on}]
How many pen lifts does \zh{讠} take? (\textbf{One}.)

Source: \href{https://www.unicode.org/charts/PDF/U4E00.pdf}{Unicode CJK Unified Ideographs chart}; stroke order per the PRC standard \emph{Xiandai Hanyu Tongyongzi Bishun Guifan} (1997).
\end{tcolorbox}
\section[xiè]{\zh{谢} — twelve strokes, and three pieces}

\label{lesson:ZH-W10-xie}



\begin{quote}
Write \zh{讠}. Dot, then one turning stroke.
\end{quote}

\subsection*{Script — the shape}
\begin{quote}
\zh{谢}
\end{quote}

\textbf{Twelve strokes.} That is the longest thing in this book, and the number is about to stop mattering.

You were told once that stroke-counting gives out and structure takes over. This is where that pays off. \textbf{\zh{谢} is three pieces, left to right}, and you have just written the first:

\begin{quote}
\textbf{\zh{讠}} + \textbf{\zh{身}} + \textbf{\zh{寸}}
\end{quote}

\begin{itemize}
\raggedright
  \item \textbf{\zh{讠}} — speech, two strokes, done
  \item \textbf{\zh{身}} — a shape meaning \emph{body}, seven strokes
  \item \textbf{\zh{寸}} — a shape meaning \emph{inch}, three strokes on the right
\end{itemize}

\zh{身} and \zh{寸} are pieces, not characters you are being asked to own. Read them, place them, move on.

\textbf{\zh{讠} (strokes 1-2).} As you just practised.

\textbf{\zh{身} (strokes 3-9).} An upper falling stroke; the left side down; the top and right side in one turning stroke that hooks left at the base; two short inner horizontals; a wide horizontal beneath them; then a falling stroke down-left.

\textbf{\zh{寸} (strokes 10-12).} A horizontal; a vertical that hooks left at the base; and a dot, last.

Eleven pen lifts. \textbf{Left piece finished before the middle, middle before the right} — the oldest habit you have, applied to three pieces instead of two.

Compare it honestly with \zh{我}. \zh{我} was seven strokes and had \textbf{no} pieces you could hold on to. \zh{谢} is twelve and has three. \textbf{The long one is the easier one.} That is the whole argument for learning structure rather than counting.

\subsection*{Your turn}
\begin{itemize}
\raggedright
  \item \emph{Write it:} \zh{讠}, then the middle piece, then the right piece — stopping between each
  \item \emph{Write it:} \zh{谢} straight through, and say \textquotedblleft{}speech, body, inch\textquotedblright{} as you go
  \item \emph{Say it:} \textbf{xiè} — falling
\end{itemize}

\begin{tcolorbox}[breakable,colback=teal!4,colframe=teal!35!black,title={Before you move on}]
Which was harder to learn, this one or \zh{我}? (\textbf{\zh{我}} — and it is five strokes shorter.)

Source: \href{https://www.unicode.org/charts/PDF/U4E00.pdf}{Unicode CJK Unified Ideographs chart}; stroke order per the PRC standard \emph{Xiandai Hanyu Tongyongzi Bishun Guifan} (1997).
\end{tcolorbox}
\section[xièxie]{\zh{谢谢} — said twice}

\label{lesson:ZH-C10-xiexie}



\begin{quote}
Say \textbf{xiè} — falling, from high to low.
\end{quote}

\begin{cousinweb}
\textbf{\zh{谢谢}} \emph{xièxie} — \textbf{thank you}.

It is \textbf{\zh{谢} written twice}. Nothing else. You learned one character and you have the word.

This is new, and it is not a quirk. Mandarin builds a great many words by \textbf{saying a piece twice}, and doubling usually softens something — it makes an action smaller, lighter, more everyday. That pattern is worth knowing; you will meet it again.

Two cautions about this particular word, so you do not over-apply it. \textbf{\zh{谢谢} is a fixed expression}, not something you assemble fresh each time — it does not mean \textquotedblleft{}to thank briefly\textquotedblright{}. And \textbf{\zh{谢} is not said on its own}: it lives inside words, so do not reach for a bare \zh{谢} and expect it to work.
\end{cousinweb}

\begin{sounds}
\emph{xièxie}, not \emph{xièxiè}.

The first \zh{谢} falls. The second \textbf{goes neutral} — short, light, no tone mark. You have seen this twice already, on \zh{上} inside \zh{早上} and on \zh{字} inside \zh{名字}, and now on a syllable that is the \emph{same character} as the one before it.

That is the useful lesson: \textbf{the light syllable is not a property of the character, it is a property of the word.} \zh{谢} falls in this word's first half and goes light in its second — and you will meet \zh{谢} falling again, in its second position, before this chapter is out. {[}REPEAT x2{]}

Say it: \emph{xiè} dropping, then \emph{xie} small and quick behind it.
\end{sounds}

\subsection*{Your turn}
\begin{itemize}
\raggedright
  \item \emph{Say it:} \textbf{xiè} alone, falling — then \textbf{xièxie}, with the second syllable light
  \item \emph{Say it:} \textbf{xièxie} twice, and say which of the two \zh{谢} carries a tone
  \item \emph{Write it:} \zh{谢谢} — the same twelve strokes, twice
\end{itemize}

\begin{tcolorbox}[breakable,colback=teal!4,colframe=teal!35!black,title={Before you move on}]
How many new characters did this word cost you? (\textbf{One}.)

Source: \href{https://www.unicode.org/charts/PDF/U4E00.pdf}{Unicode CJK Unified Ideographs chart}
\end{tcolorbox}
\section[bú xiè]{\zh{不谢} — answering thanks}

\label{lesson:ZH-C10-buxie}



\begin{quote}
Say \textbf{xièxie} — falling, then light.
\end{quote}

\begin{cousinweb}
Someone thanks you. What do you say back?

\begin{quote}
\zh{不谢}
\end{quote}

\textbf{\zh{不谢}} — and it is made of two things you already have. \textbf{\zh{不}}, the negator you have carried since \zh{好} and \zh{不好}, in front of \textbf{\zh{谢}}, which you learned an hour ago. Nothing new. You could have built this yourself.

It works the way \zh{不是} works: \textbf{put \zh{不} in front of the word and you turn it.} Here turning it is a polite refusal of the thanks — \emph{no need to thank me} — which is how a great many languages answer being thanked.
\end{cousinweb}

\begin{sounds}
Before you say it, work it out. \textbf{\zh{谢} is fourth tone.} What does \zh{不} do in front of a fourth tone?

\textbf{It rises.} \emph{bú xiè}, not \emph{bù xiè} — exactly as \zh{不是} became \emph{bú shì}.

This is the second time that rule has fired, and this time you were told the condition before the word. That is the difference between a rule you have been given and a rule you can use: \textbf{you can now predict the pronunciation of a word nobody has said to you yet.} {[}REPEAT x2{]}
\end{sounds}

\begin{culture}
\zh{不谢} is real and correct, and it is not the commonest reply. The one you will hear most often needs two characters this book cannot yet print, so it is waiting. \zh{不谢} will be understood everywhere, and it has the advantage of costing you nothing you did not already have.
\end{culture}

\subsection*{Your turn}
\begin{itemize}
\raggedright
  \item \emph{Say it:} \textbf{bú xiè} — the first syllable rising, and know why
  \item \emph{Say it:} \textbf{bú shì}, then \textbf{bú xiè} — the same rule, twice
  \item \emph{Write it:} \zh{不谢}
\end{itemize}

\begin{tcolorbox}[breakable,colback=teal!4,colframe=teal!35!black,title={Before you move on}]
How many new characters did this reply cost you? (\textbf{None}.)

Source: \href{https://www.unicode.org/charts/PDF/U4E00.pdf}{Unicode CJK Unified Ideographs chart}
\end{tcolorbox}
\section[Practice]{Practice — thanks, and the answer}

\label{lesson:ZH-C10-practice}



\begin{quote}
Say \textbf{bú xiè}. Rising, then falling — and you know why the first one rises.
\end{quote}

\subsection*{The exchange}
\begin{quote}
\textbf{A:} \zh{早上好}
\end{quote}

>

\begin{quote}
\emph{zǎoshang hǎo} — good morning.
\end{quote}

>

\begin{quote}
\textbf{B:} \zh{早上好}
\end{quote}

>

\begin{quote}
\emph{zǎoshang hǎo} — good morning.
\end{quote}

>

\begin{quote}
\textbf{A:} \zh{谢谢}
\end{quote}

>

\begin{quote}
\emph{xièxie} — thank you.
\end{quote}

>

\begin{quote}
\textbf{B:} \zh{不谢}
\end{quote}

>

\begin{quote}
\emph{bú xiè} — you're welcome.
\end{quote}

>

\begin{quote}
\textbf{A:} \zh{再见}
\end{quote}

>

\begin{quote}
\emph{zàijiàn} — goodbye.
\end{quote}

Six lines, and \textbf{one new character} paid for two of them.

\begin{grammarlens}[title={What you've built — one character, a whole exchange}]
This chapter cost you \zh{谢} and nothing else. Out of it came the thanks (\textbf{\zh{谢谢}}, the same character twice) and the reply (\textbf{\zh{不谢}}, a negator you already had). That is what a system does for you that a phrasebook cannot: pieces recombine.

Three things you did here without being led:

\begin{itemize}
\raggedright
  \item you \textbf{doubled} a character to make a word, and knew the second half would go light;
  \item you \textbf{negated} a word to make its answer, the same move that turned \zh{是} into \zh{不是};
  \item and you \textbf{predicted a tone change} from a rule, before hearing the word said.
\end{itemize}

That last one is the real marker. Early on you were told what \zh{你好} sounds like. Now you are told a rule and you work out what a word sounds like. Those are different kinds of knowing, and only the second one travels to words nobody has taught you.
\end{grammarlens}

\subsection*{Your turn}
\begin{itemize}
\raggedright
  \item \emph{Say it:} the whole exchange, both parts
  \item \emph{Say it:} \textbf{xièxie}, then \textbf{bú xiè} — light on the second, rising on the first
  \item \emph{Write it:} \zh{谢谢}, then \zh{不谢}
\end{itemize}

\begin{tcolorbox}[breakable,colback=teal!4,colframe=teal!35!black,title={Before you move on}]
Which two of these did you work out rather than being told? (\textbf{The light second syllable, and the rising \zh{不}}.)

Source: \href{https://www.unicode.org/charts/PDF/U4E00.pdf}{Unicode CJK Unified Ideographs chart}
\end{tcolorbox}
